\documentclass{article}

\usepackage{amsmath,amssymb}
\usepackage{xurl}
\usepackage[hidelinks]{hyperref}

% Shared commands for notes in docs/paper-gaps/.

\DeclareUnicodeCharacter{2080}{\ifmmode{}_{0}\else\textsubscript{0}\fi}
\DeclareUnicodeCharacter{2081}{\ifmmode{}_{1}\else\textsubscript{1}\fi}
\DeclareUnicodeCharacter{2082}{\ifmmode{}_{2}\else\textsubscript{2}\fi}
\DeclareUnicodeCharacter{2099}{\ifmmode{}_{n}\else\textsubscript{n}\fi}

\newcommand{\Lean}{\textsc{Lean}}
\ExplSyntaxOn
\NewDocumentCommand{\leanid}{m}{
  \ifmmode
    \textsf{\detokenize{#1}}
  \else
    \group_begin:
    \urlstyle{sf}
    \tl_set:Nn \l_tmpa_tl {#1}
    \tl_replace_all:Nnn \l_tmpa_tl {\_} {_}
    \exp_args:NV \path \l_tmpa_tl
    \group_end:
  \fi
}
\ExplSyntaxOff
\providecommand{\tr}{\operatorname{tr}}
\newcommand{\tnleanIssueBase}{https://github.com/LionSR/TNLean/issues/}
\newcommand{\tnleanPrBase}{https://github.com/LionSR/TNLean/pull/}
\newcommand{\tnleanDocBase}{https://sirui-lu.com/TNLean/docs/}
\newcommand{\ghissue}[1]{\href{\tnleanIssueBase#1}{GitHub issue \##1}}
\newcommand{\ghpr}[1]{\href{\tnleanPrBase#1}{PR \##1}}
\newcommand{\doclink}[1]{\href{\tnleanDocBase#1.html}{\path{#1}}}

% Dirac notation and the identity operator, as in blueprint/src/macros/common.tex.
% \providecommand keeps note-local definitions authoritative.
\usepackage{dsfont}
\providecommand{\ket}[1]{|#1\rangle}
\providecommand{\bra}[1]{\langle #1|}
\providecommand{\braket}[2]{\langle #1 | #2 \rangle}
\providecommand{\ketbra}[2]{|#1\rangle\!\langle #2|}
\providecommand{\Id}{\mathds{1}}
\providecommand{\one}{\mathds{1}}
\providecommand{\C}{\mathbb{C}}
\providecommand{\MN}[1]{M_{#1}(\C)}
\providecommand{\MD}{M_D(\C)}

% Verdict marker: \gapnote{<kind>}{<status>}. Kind is the policy
% classification (clarification, local-correction, scope-restriction,
% unfaithful, false-source, open-gap); status is open, wip (elimination
% underway), resolved, or historical. Machine-read by the site index;
% prints a verdict line.
\newcommand{\gapnote}[2]{\par\noindent\textbf{Verdict:} #1 (#2).\par}


\title{Blocking in the Symmetry-Preserving Standard-Form Path Criterion}
\author{}
\date{2026-11-16}

\begin{document}
\maketitle
\gapnote{scope-restriction}{open}

\section{The printed lemma}

Lemma~\texttt{lem:equivalent-symmetry} of Ref.~\cite{Cirac2017MPU}
(lines~1369--1376) states that two MPUs $U_{1,2}^{(N)}$ described by simple
tensors are strictly equivalent under a symmetry $\mathcal S$ if and only if
there are continuous families of unitaries $\tilde u(p),\tilde v(p)$,
$p\in[0,1]$, whose endpoints generate $U_1^{(N)}$ and $U_2^{(N)}$ and such that
the MPU generated by $\tilde u(p),\tilde v(p)$ is invariant under $\mathcal S$
for every $p$. The proof reads ``The proof follows directly from Corollary
\texttt{coruvSF}'' (line~1379). Strict equivalence under a symmetry is
Definition~\texttt{def:strictly-equivalent-symmetry} (lines~1345--1354): a
continuous path of tensors, not necessarily in canonical form, joining the two
tensors through MPUs invariant under $\mathcal S$.

\section{The two directions}

The MPU generated by gates $u,v$ is the operator on the chain of standard-form
sites. Following the conventions recorded in
\url{https://sirui-lu.com/TNLean/paper-gaps/mpu_standard_form_canonical_convention.pdf}
and
\url{https://sirui-lu.com/TNLean/paper-gaps/mpu_standard_form_parity_gap.pdf},
a tensor in standard form is the two-site block of a simple tensor in
canonical form~II, and $U^{(N)}$ counts two-site blocks.

\paragraph{Reverse direction.} Let the gates $\tilde u(p),\tilde v(p)$ act on
blocks of $k$ sites, with endpoints the standard-form gates of the $k$-site
blocks of $\mathcal U_1$ and $\mathcal U_2$. The tensors built from the gates
form a continuous path of tensors on blocks of $2k$ sites, and interpolating
half factors joins each endpoint to the tensor of its own gates without
changing any periodic operator. So the blocks of $2k$ sites of
$\mathcal U_1$ and $\mathcal U_2$ are strictly equivalent under the blocked
symmetry. For $k=1$ this is the printed reverse direction, with lengths
counted on the standard-form chain.

\paragraph{Forward direction.} A strict equivalence gives a path
$\mathcal W(p)$ of MPU tensors, but the tensors along it need not be simple.
Corollary \texttt{coruvSF}, which the proof invokes, produces continuous gates
only for the MPUs $U^{(2k_0N)}(p)$ with $k_0\le D^4$, where $D$ is the bond
dimension of the path. The formalized forward direction therefore holds after
blocking $2k_0$ sites: the endpoint gates are those of the standard forms of
the $k_0$-site blocks of $\mathcal U_1$ and $\mathcal U_2$, and invariance
is imposed at the lengths $2k_0N$.

\section{Scope}

The printed lemma is used with the same $U^{(N)}$ in both directions. The
formalization proves the reverse direction for every blocking length, in
particular on the standard-form chain, and the forward direction after
blocking $2k_0$ sites. Whether the forward direction holds without the extra
blocking, that is, whether a strict equivalence of simple tensors under
$\mathcal S$ yields gates for the standard forms of the tensors themselves,
is open. The distinction matters because strict equivalence and equivalence
after blocking differ in the presence of a symmetry, as the source notes at
line~1366.

Both directions are stated for symmetries imposed at positive chain lengths.
The periodic operator at length zero records only the bond dimension, which
the source does not consider.

\section{Formal statements}

The reverse direction is
\leanid{MPOTensor.strictlyEquivalentUnderSymmetry_blockTensor_of_gate_path},
with the case $k=1$
\leanid{MPOTensor.strictlyEquivalentUnderSymmetry_blockTwo_of_gate_path}.
The blocked forward direction is
\leanid{MPOTensor.exists_gate_path_of_strictlyEquivalentUnderSymmetry}, which
feeds Corollary \texttt{coruvSF} into
\leanid{MPOTensor.gatePath_of_symmetricPath}; its gates act on blocks of some
positive number $k$ of sites, and their generated MPUs are invariant under
the symmetry blocked $2k$ sites. Lemma~\texttt{lem:equivalent-symmetry} itself remains unformalized.

\bibliographystyle{alpha}
\bibliography{references}
\end{document}
